\section{Experimental Setup}
\label{sec:experimental-setup}

\subsection{Datasets}

We evaluate WR-STV using three trace sets collected from the
TrainTicket benchmark microservice system~\cite{zhou2021trainticket}.
Each trace set corresponds to a distinct configuration of
\texttt{ts-order-service}:
\texttt{ts-order-service\_mongodb\_4.2.2},
\texttt{ts-order-service\_mongodb\_4.4.15}, and
\texttt{ts-order-service\_3.0.4-mongodb-driver}. We refer to these
configurations as Dataset~1, Dataset~2, and Dataset~3, respectively.
TrainTicket is a benchmark microservice system designed for fault
analysis and debugging studies, making it suitable for controlled
experiments on microservice observability and anomaly scoring.

The three configurations vary the service version and MongoDB driver
while preserving broadly similar service-level functionality. This makes
them useful for both within-dataset evaluation and cross-dataset transfer
analysis. In particular, the related but non-identical configurations
allow us to study how WR-STV behaves when training and test traces share
similar service structures but differ in observed latency distributions
and service-path coverage.

Table~\ref{tab:dataset-stats} summarizes trace, span, and service
counts, along with trace-graph quality checks performed before modeling.
Across all three datasets, each trace contains exactly one root span, no
non-root span has a missing parent, and no duplicate
\((\text{traceID}, \text{spanID})\) pairs are present.

\begin{table}[t]
\centering
\caption{Dataset statistics and trace-graph quality checks.}
\label{tab:dataset-stats}
\begin{tabular}{lccc}
\toprule
& \textbf{Dataset 1} & \textbf{Dataset 2} & \textbf{Dataset 3} \\
\midrule
Traces          & 269 & 86  & 160 \\
Spans           & 3{,}788 & 1{,}306 & 1{,}993 \\
Distinct services & 32 & 33 & 33 \\
\midrule
Exactly 1 root span/trace & \checkmark & \checkmark & \checkmark \\
Non-root spans w/ missing parent & 0 & 0 & 0 \\
Duplicate (traceID, spanID) pairs & 0 & 0 & 0 \\
\bottomrule
\end{tabular}
\end{table}


\subsection{Preprocessing Pipeline}

Spans are grouped by trace identifier, and each span's root-to-node
service path is reconstructed by following parent-child relationships
from the trace root. Following Equation~\ref{eq:stv-value}, spans that
share the same \((\text{service}, \text{service\_path})\) feature within
a single trace are aggregated into one STV feature by retaining the
maximum observed duration. We use maximum aggregation because the slowest
occurrence of a repeated service path is often the most actionable signal
for latency anomaly triage, whereas mean aggregation may dilute localized
delays.

For every experimental split, the STV vocabulary, feature-wise
normalization statistics, and reliability weights are estimated only from
the training traces. The learned vocabulary and statistics are then
applied unchanged to validation and test traces. This protocol prevents
test-set leakage and ensures that unseen service paths at inference time
are treated as absent rather than being used to expand the feature space.

\subsection{Synthetic Latency Injection}

Because the trace sets do not provide per-trace ground-truth anomaly
labels, we evaluate anomaly scoring using controlled synthetic latency
injection. For a normal test trace, an injected copy is created by
multiplying the duration of selected STV-visible spans by an injection
factor \(f\). A span is considered STV-visible if its
\((\text{service}, \text{service\_path})\) feature appears in the
training vocabulary. The original trace is assigned label \(0\), while
the injected copy is assigned label \(1\).

We evaluate two span-selection strategies:

\begin{itemize}
    \item \textbf{Largest-visible injection:} selects up to \(m\)
    STV-visible spans with the largest original durations. This is the
    primary controlled latency-injection setting.

    \item \textbf{Random-visible injection:} selects up to \(m\)
    STV-visible spans uniformly at random, regardless of original
    duration. This setting tests whether WR-STV's advantage depends on
    always perturbing high-duration spans.
\end{itemize}

Unless otherwise specified, experiments use \(f=30\) and \(m=3\). We
also perform injection-strength sensitivity analysis using
\(f \in \{5,10,20,30\}\), as reported in
Section~\ref{sec:results-sensitivity}. Both injection strategies are
restricted to STV-visible dimensions, meaning that they evaluate
controlled latency deviations on service paths already observed during
training. This is a deliberate scope limitation and is discussed further
in Section~\ref{sec:threats}.

\subsection{Baseline Methods}

We compare WR-STV against both raw-STV and residual-STV baselines:

\begin{itemize}
    \item \textbf{Baseline STV + Isolation Forest:} the raw STV
    representation from Equation~\ref{eq:stv-value}, without feature
    normalization, scored using Isolation Forest~\cite{liu2008iforest}.
    This is the primary baseline used to measure trace-complexity bias.

    \item \textbf{Enhanced STV + Isolation Forest:} z-score-normalized
    residual features from Equations~\ref{eq:mean}--\ref{eq:residual},
    scored using Isolation Forest. This baseline isolates the effect of
    residual normalization independent of WR-STV's top-\(k\) scoring
    rule and reliability weighting.

    \item \textbf{Max-Z Residual STV:} the maximum absolute residual from
    Equation~\ref{eq:residual}, using \(k=1\) and no reliability
    weighting. This isolates the contribution of the reliability weight
    in Equation~\ref{eq:weight}.

    \item \textbf{Top-3 Residual STV:} the mean of the top three
    unweighted absolute residuals. This tests whether aggregating
    multiple residual contributors changes ranking performance without
    reliability weighting.

    \item \textbf{Residual STV + LOF / One-Class SVM:} the standardized
    residual representation scored using Local Outlier Factor
    \cite{breunig2000lof} and One-Class SVM~\cite{scholkopf2001ocsvm}.
    This compares WR-STV's simple top-\(k\) aggregation against classical
    detectors applied to the same residual feature space.

    \item \textbf{Raw STV + LOF / One-Class SVM:} the same two detectors
    applied directly to the unnormalized STV representation. This tests
    whether the behavior observed with Isolation Forest is specific to
    one detector family or also appears under other classical
    unsupervised methods.
\end{itemize}

\subsection{Evaluation Metrics}

We report ROC-AUC, average precision (AUPRC), F1 at an oracle threshold,
precision@\(N_+\), and pairwise win rate. Here, \(N_+\) denotes the
number of injected traces in an evaluation split. The oracle-threshold
F1 is computed by labeling the top \(N_+\) scored traces as anomalous and
is therefore used only as a ranking diagnostic, not as a deployment-ready
thresholding result.

Pairwise win rate is computed over matched
\((\text{original}, \text{injected})\) trace pairs. A pair is counted as a
win when the injected trace receives a strictly higher anomaly score than
its corresponding original trace. This metric directly measures whether
a scoring method ranks a controlled latency perturbation above the normal
trace from which it was generated.

\subsection{Hyperparameter Selection}
\label{sec:hyperparameter-selection}

The WR-STV top-\(k\) parameter and residual clip bound \(C\) are selected
using an inner validation split. We search over
\(k \in \{1,3,5\}\) and \(C \in \{3.0,5.0,10.0\}\). The validation split
is drawn from the outer training set, with \(30\%\) of the outer training
traces reserved for validation. Test traces are not used during
hyperparameter selection. This procedure selects \(k=1\) and \(C=5.0\),
which are used for all final WR-STV results.

\begin{table}[t]
\centering
\caption{Hyperparameter selection on a held-out validation split
(independent of all test-set metrics reported elsewhere in this paper).}
\label{tab:hyperparam-validation}
\begin{tabular}{lc}
\toprule
\textbf{Hyperparameter} & \textbf{Selected value} \\
\midrule
$k$ (top-$k$ features aggregated) & 1 \\
$C$ (residual clip bound)         & 5.0 \\
\bottomrule
\end{tabular}

\vspace{4pt}
\footnotesize
\end{table}


\subsection{Implementation Details}

All experiments are implemented in Python using NumPy, pandas, and
scikit-learn. Isolation Forest uses 200 estimators in all experiments.
Local Outlier Factor is evaluated in novelty-detection mode, and
One-Class SVM uses the radial-basis-function kernel with
\(\nu=0.1\) and \(\gamma=\texttt{scale}\). Unless otherwise specified,
each dataset is split 65:35 into training and test traces. The split is performed uniformly at random at the trace level using the specified random seed, ensuring that individual traces appear in only one partition. Within-dataset
experiments are repeated over five random seeds for each dataset,
resulting in \(3 \times 5 = 15\) runs.

The overall experimental protocol is as follows. First, STVs are
constructed from the training split, and the training vocabulary is fixed.
Second, WR-STV feature statistics and reliability weights are estimated
from the same training split only. Third, the fitted scoring procedure is
applied unchanged to held-out validation or test traces. Finally,
synthetic latency anomalies are injected only into evaluation traces, and
metrics are computed over the resulting normal-vs-injected trace pairs.